\section{Discussão}
\label{sc:discussao}

\subsection{Qual método escolher}

Em relação à QP1, nenhum método se destaca estatisticamente pela similaridade de cosseno: as diferenças entre os seis são pequenas diante da variação entre perguntas, e a ordenação muda de domínio para domínio. Pela \emph{Answer Relevancy} há diferenças significativas, e em custo as diferenças são grandes e estáveis. Nessas condições, as recomendações mais sólidas são de exclusão. O Reciprocal custa cerca de cinco vezes o Stuff em tokens, tem a pior \emph{Answer Relevancy} e fica no grupo intermediário em similaridade: no conjunto das perguntas, não compensa o custo adicional. O Query Step-Down teve a maior \emph{Answer Relevancy}, mas sem diferença significativa em relação ao Map-Reduce, que usa cerca de um terço dos tokens (1.987 contra 6.773 por pergunta).

Entre os métodos restantes, a escolha depende do que se entende por qualidade. O Map-Reduce foi o mais equilibrado: segundo em \emph{Answer Relevancy}, intermediário em similaridade e com custo de 1,5 vez o Stuff. O Map-Rerank, com o mesmo custo, produz respostas curtas e diretas, o que lhe deu a maior similaridade média e uma das piores \emph{Answer Relevancy}. Ele é adequado quando a resposta esperada é um valor único, como nas perguntas financeiras, e menos quando se espera uma resposta explicativa. O Refine teve a maior \emph{Faithfulness} bruta, sem diferença significativa, a um custo de 1,9 vez o Stuff. O Stuff continua sendo a opção mais barata e mais rápida, com a ressalva das recusas, discutida a seguir.

Esses resultados divergem em parte do estudo de referência \cite{sakar2025rag}, no qual o Reciprocal RAG obteve a maior similaridade. Aqui ele foi o melhor apenas no domínio financeiro. Três diferenças de desenho podem explicar a divergência: o idioma, o tipo de pergunta (questões discursivas e perguntas factuais sobre documentos longos) e a escala, já que 27 perguntas não permitem detectar diferenças pequenas entre métodos.

\subsection{Recusas explicam boa parte da diferença do Stuff}

A análise das recusas sugere um mecanismo concreto por trás da desvantagem do Stuff. Ele recebe os quatro fragmentos de uma vez, com a instrução de dizer que não sabe quando a resposta não estiver no contexto, e recusou 30\% das perguntas. Excluídas as recusas, sua similaridade fica no mesmo patamar da do Map-Rerank, e sua \emph{Faithfulness} passa a ser a maior entre os métodos. O Map-Rerank, ao avaliar cada fragmento separadamente e reter a resposta de maior confiança, é pouco afetado por recusas: basta um fragmento que sustente uma resposta. As recusas se concentraram no domínio médico, cujas perguntas combinam uma vinheta clínica longa com um enunciado específico e cujo corpus reúne cinco documentos.

Barnett et al.~\cite{barnett2024seven} distinguem dois pontos de falha relevantes aqui: o contexto correto não é recuperado, ou é recuperado e o modelo não consegue extrair dele a resposta. As métricas RAGAS indicam que o segundo caso predominou. Em 10 das 13 recusas, o juiz considerou que o contexto recuperado cobria ao menos metade da resposta de referência, com a ressalva de que o próprio \emph{Context Recall} é uma estimativa ruidosa (Seção~\ref{ssc:disc_juiz}). As recusas também são instáveis: o Stuff foi executado três vezes sobre as mesmas perguntas, com temperatura 0, e em uma delas recusou em uma execução e respondeu corretamente nas outras duas (similaridade de 0,433 contra 0,87). Parte da diferença entre métodos, portanto, pode depender do texto do prompt e de decisões do modelo que estão no limite, e não só da arquitetura. Uma instrução de recusa mais branda aproximaria o Stuff dos demais métodos, possivelmente ao custo de mais respostas sem sustentação no contexto. Isso não foi testado aqui e fica como trabalho futuro.

\subsection{A métrica muda a conclusão}

A resposta à QP2 é afirmativa: a ordenação dos métodos depende da métrica. O Map-Rerank passa do primeiro lugar em similaridade de cosseno para o quinto em \emph{Answer Relevancy}, e o Reciprocal, intermediário no cosseno, fica em último. As correlações linha a linha entre cosseno e as métricas RAGAS são fracas ($\rho$ entre 0,16 e 0,35). O mecanismo é a forma da resposta. O cosseno premia respostas curtas quando a referência também é curta, e penaliza respostas corretas redigidas em frase completa. A \emph{Answer Relevancy} faz o oposto. Nenhuma das duas verifica o valor em si: na pergunta sobre o patrimônio líquido, um número diferente do gabarito obteve similaridade de 0,613.

Isso tem uma consequência para estudos comparativos como o de referência, que usam apenas a similaridade de cosseno: parte da vantagem atribuída a um método pode refletir a forma das respostas que ele produz, e não sua correção. A recomendação que decorre deste trabalho é reportar mais de uma métrica e declarar qual noção de qualidade importa para a aplicação: exatidão de um valor, completude da explicação ou ancoragem no documento. Para respostas numéricas, uma comparação direta do valor extraído seria mais adequada que qualquer métrica baseada em embeddings.

\subsection{Confiabilidade do juiz automático}
\label{ssc:disc_juiz}

A validação do juiz não indicou viés de autopreferência. Nas métricas que avaliam a própria resposta, o \texttt{gpt-4o-mini} foi igualmente ou mais rigoroso com as respostas que gerou do que um juiz de outra família. Isso apoia o uso do mesmo modelo como gerador e juiz neste estudo, como também fizeram Medeiros e Oliveira~\cite{medeiros2025embeddings}, ao menos para as métricas e o modelo avaliados aqui.

A confiabilidade, porém, depende da métrica, e isso determina o que pode ser concluído de cada uma. A \emph{Answer Relevancy} foi estável no reteste e entre juízes, com a mesma ordenação dos métodos. As conclusões baseadas nela, inclusive a inversão em relação ao cosseno discutida acima, são robustas à escolha do juiz. O \emph{Context Recall} também foi estável, o que sustenta a leitura das recusas. A \emph{Faithfulness} variou de forma considerável mesmo quando o mesmo juiz avaliou duas vezes a mesma resposta, e os juízes divergiram no tratamento das recusas. As diferenças de \emph{Faithfulness} entre métodos, que já não eram significativas, não devem ser interpretadas além disso. A \emph{Context Precision} não foi confiável: os juízes não concordaram entre si, e a mesma entrada recebeu notas diferentes em avaliações repetidas. Como ela avalia a recuperação, que é idêntica entre os métodos, isso não afeta a comparação entre eles, mas seus valores absolutos não devem ser tomados como medida da qualidade da recuperação.

Do ponto de vista metodológico, a validação custou pouco: 54 respostas reavaliadas por dois juízes. Estudos que usam LLMs como juízes com um único modelo e uma única execução por resposta tratam como medida o que, em parte, é ruído do juiz. Um reteste numa amostra e a comparação com um juiz independente indicam quais métricas sustentam conclusões e quais não.

\subsection{Compressão contextual}

Em relação à QP3, os limiares fixos não se mostraram transferíveis. Com o BGE-M3, a similaridade dos fragmentos recuperados ficou concentrada numa faixa estreita, entre 0,52 e 0,75. Limiares abaixo dela, como 0,3 e 0,5, não têm efeito; limiares acima, como 0,7, reduzem o contexto a um único fragmento. Valores definidos para outros modelos de embedding, como os do estudo de referência, não se aplicam diretamente, porque a escala de similaridade é própria de cada modelo. A calibração pelos quantis da distribuição observada resolve esse problema: transforma o limiar em uma fração aproximada de fragmentos descartados e produz um compromisso gradual entre tokens e qualidade.

Os valores calibrados foram parecidos nos três domínios (quantil 50 entre 0,626 e 0,639), o que sugere que, para um mesmo embedding, a distribuição de similaridade é relativamente estável. O que variou foi o quanto cada domínio depende dos fragmentos descartados. Nos domínios financeiro e técnico-científico, com um único documento cada, descartar metade dos fragmentos economizou mais de 40\% dos tokens sem perda de similaridade. No domínio médico, em que a resposta pode depender de fragmentos de documentos diferentes e as perguntas são longas, cada fragmento descartado custou qualidade. A recomendação prática é calibrar o limiar em uma amostra de perguntas do domínio de interesse e validar a qualidade antes de adotá-lo. Como aqui a calibração usou as mesmas perguntas avaliadas, os ganhos observados são um limite superior do que se obteria com um limiar escolhido de antemão.

\subsection{Bancos vetoriais}

Em relação à QP4, a comparação confirma que, com busca exata e a mesma métrica de distância, o banco vetorial não altera a qualidade das respostas: os fragmentos recuperados foram idênticos. O FAISS foi muito mais rápido e econômico em memória, mas os tempos do ChromaDB, de poucos milissegundos, são desprezíveis diante dos segundos gastos na geração. Para corpora do tamanho dos usados aqui, de dezenas a centenas de fragmentos, a persistência e os filtros por metadados do ChromaDB custam pouco. A vantagem do FAISS tende a importar em corpora ordens de grandeza maiores ou com alto volume de consultas, cenários que não foram avaliados.

\subsection{Limitações}

Os resultados devem ser lidos à luz das seguintes limitações.

\begin{itemize}
  \item \textbf{Tamanho da amostra.} São 27 perguntas no total e de 7 a 12 por domínio. Isso basta para medir diferenças de custo, que são grandes, mas não para detectar diferenças moderadas de qualidade entre métodos.
  \item \textbf{Configuração fixa.} Um único LLM (\texttt{gpt-4o-mini}), um único embedding (BGE-M3), um único tamanho de fragmento e $k = 4$. As conclusões valem para essa configuração. Com um LLM mais capaz, por exemplo, a taxa de recusas e a vantagem do Map-Rerank poderiam ser diferentes.
  \item \textbf{Perguntas redigidas pelos autores.} Nos domínios financeiro e técnico-científico, as perguntas e as respostas de referência foram redigidas pelos autores, ainda que ancoradas em trechos literais. Só o domínio médico tem gabarito oficial externo.
  \item \textbf{Extração de texto.} Os PDFs foram convertidos em texto simples, e tabelas financeiras perdem parte da estrutura nessa conversão, o que afeta a recuperação de valores numéricos.
  \item \textbf{Métricas automáticas.} A similaridade de cosseno depende da forma da resposta, e as métricas RAGAS dependem de um juiz LLM não determinístico (Seção~\ref{ssc:disc_juiz}). Nenhuma das duas foi validada contra avaliação humana.
  \item \textbf{Determinismo da geração.} Mesmo com temperatura 0, a geração não é totalmente determinística. O Stuff sem filtro foi executado três vezes sobre as 27 perguntas (no comparativo e nas duas rodadas de compressão), e 21 respostas saíram idênticas nas três execuções. Cada resultado reportado vem de uma única execução por pergunta e método.
  \item \textbf{Custo de hardware.} Como a geração ocorre em uma API remota, as medidas locais de CPU e memória não refletem o custo de inferência dos métodos. O custo foi medido em tokens e tempo, e a comparação de hardware se restringiu aos bancos vetoriais.
  \item \textbf{Versão do modelo.} Os registros armazenaram o apelido do modelo, e não a versão exata, o que dificulta a reprodução exata caso o apelido passe a apontar para outra versão.
\end{itemize}
